\chapter{Number of Planets}
\section{chatGPT}

When I studied in the 1970s, Pluto was officially counted as the ninth planet. The change happened in 2006, when astronomers adopted a more precise definition of what counts as a planet.
What changed?
The problem arose because astronomers began discovering many objects beyond Neptune, in the Kuiper Belt, some of them surprisingly large. In particular, Eris, discovered in 2005, is nearly as large as Pluto.
This raised a difficult question:
If Pluto is a planet, should Eris also be a planet? And what about the other large Kuiper Belt objects that might subsequently be discovered?
So in 2006 the International Astronomical Union (IAU) adopted a definition requiring a body to satisfy three conditions to be a planet in our Solar System:
\begin{enumerate}
   \item It orbits the Sun.
   \item It is massive enough for its own gravity to make it approximately round.
   \item It has cleared the neighborhood around its orbit of other objects.
\end{enumerate}
Pluto satisfies (1) and (2), but not (3). It lives in the Kuiper Belt among many other objects of comparable orbital character.
Consequently, Pluto was reclassified as a dwarf planet, leaving eight planets:
Mercury – Venus – Earth – Mars – Jupiter – Saturn – Uranus – Neptune.
\paragraph{The interesting part}
Nothing physically happened to Pluto in 2006!
It did not become smaller, change its orbit, or somehow "lose" planetary properties. The change was taxonomic: astronomers changed the definition of the word planet.

In fact, there is a rather amusing historical twist. Pluto was discovered in 1930, and its status as a planet was already somewhat unusual because it is much smaller than the other planets. For decades this was not a serious problem because no comparable objects were known.
Once the Kuiper Belt began to reveal itself, Pluto turned out to be the first known member of a much larger population.

\paragraph{Why "cleared its orbit"?}
This is actually the most controversial part of the definition.
Earth, for example, gravitationally dominates its orbital region. Pluto does not: it shares its broad neighborhood with many Kuiper Belt objects. So the distinction is roughly between:
\emph{planet} which is gravitationally dominant body in its orbital zone; and
\emph{dwarf planet} which is round body that orbits the Sun but is not dynamically dominant.

This is why Ceres, which was once classified as a planet in the 19th century, is also now a dwarf planet.
And there is a nice historical symmetry here: Ceres was considered the fifth planet for several decades, until astronomers discovered that it was one of many objects in the asteroid belt. Pluto underwent essentially the same conceptual fate more than a century later.

So old university textbooks were not "wrong" for their time. The astronomical classification changed because observations changed our understanding of the Solar System.

There is still a legitimate scientific debate about whether the 2006 definition is the best one—particularly because "clearing the neighborhood" is a dynamical criterion rather than an intrinsic physical property.
There is not universal agreement among astronomers about the 2006 definition. There is, however, an important distinction between official classification and scientific agreement.
The International Astronomical Union (IAU) officially adopted the eight-planet definition in 2006, and it has not changed that definition. Under the IAU definition, Pluto is a dwarf planet because it has not cleared its orbital neighborhood.
But the decision was controversial even at the time. The debate at the 2006 IAU meeting was quite vigorous, and roughly 400 astronomers subsequently signed a petition objecting to the adopted definition.
There are really two different concepts of "planet"
The disagreement is largely about what property should define a planet.
IAU/dynamical definition:
A planet is round, orbits the Sun, and has dynamically cleared its neighborhood.
This makes orbital dynamics part of the definition.
The alternative, advocated particularly by some planetary scientists, is a geophysical definition:
A planet is a sufficiently massive body whose gravity makes it approximately round, regardless of whether it dominates its orbital neighborhood.
Under that definition, Pluto is a planet, as are several other dwarf planets such as Ceres and Eris.
This is not merely a semantic curiosity. Consider two bodies with essentially identical physical properties. One happens to orbit alone; the other happens to be in a dense population of similar objects. Under the IAU definition, the first can be a planet and the second a dwarf planet. Critics argue that this makes planet partly a statement about the body's environment, rather than about the body itself. That criticism was already being made during the 2006 debate.
And this debate is still alive

Interestingly, this is not just an old 2006 controversy. In April 2026, Nature reported that the question had been revived publicly and described the issue as one that still divides astronomers.

And there is an amusing philosophical point here: the discovery of Pluto did not reveal that the Solar System had nine planets; the discovery of the Kuiper Belt revealed that our old concept of "planet" was inadequate. The argument is really about how we should classify a newly understood physical hierarchy.

"Clearing the orbit" is actually a rather strange criterion, 
including the quantitative dynamical parameter used to distinguish Earth/Jupiter from Pluto.
The phrase "cleared its orbit" sounds much more absolute than the underlying physics actually is. It is useful dynamically, but it is a somewhat peculiar choice for defining a category of astronomical bodies.

\paragraph{What does "cleared" actually mean?}
It does not mean that there are literally no other objects crossing the planet's orbit.
Earth, for example, has asteroids whose orbits cross or approach Earth's orbit. Jupiter has Trojan asteroids sharing its orbital period. Neptune is in a complicated dynamical relationship with many trans-Neptunian objects.
So "cleared" really means something like:
The body is gravitationally dominant in its orbital zone.
Over the age of the Solar System, a sufficiently massive body can scatter, capture, eject, or otherwise control most smaller bodies in its neighborhood.
That makes the criterion fundamentally dynamical, rather than a property intrinsic to the object.

\paragraph{The strange consequence}. The same object can be classified differently depending on its surroundings
Imagine taking Earth and putting it somewhere in the Kuiper Belt, surrounded by thousands of objects on similar orbits.
Its mass, radius, composition, gravity, geology, atmosphere, etc. would not initially change.
But it might no longer satisfy the "clearing" criterion.
Conversely, put a sufficiently massive body in a relatively empty orbital region and it can satisfy the criterion.
Thus:
$$ \text{planet status} \neq \text{property of the body alone}. $$
It depends partly on the relationship between the body and its orbital environment.
That is unusual compared with many astronomical classifications.

Compare it with the first two IAU criteria
The first two conditions are quite natural physical properties:
$$ \boxed{\text{orbits the Sun}} $$
and
$$ \boxed{\text{is massive enough to become approximately round}}. $$
The second is particularly meaningful physically. Once self-gravity dominates over material strength, a sufficiently massive body approaches hydrostatic equilibrium.
That gives us a genuine physical transition:
$$ \text{small irregular body} \quad\longrightarrow\quad \text{self-gravitating approximately spherical body}. $$
But "cleared its neighborhood" introduces a completely different kind of property:
$$ \boxed{\text{dynamical dominance over neighboring bodies}}. $$
It isn't a sharp physical transition in the same sense.
\paragraph{ There isn't a literal boundary called "the neighborhood"}
This is another problem.
How large is the neighborhood?
For a planet one might consider its Hill sphere,
$$ r_H \simeq a\left(\frac{M}{3M_\odot}\right)^{1/3}, $$
where \(a\) is the semimajor axis and \(M\) is the planet's mass.
But the Hill sphere isn't the same thing as the orbital region affected by the planet over billions of years.
You can instead define a region in terms of:
\begin{itemize}
\item mean-motion resonances,
\item scattering,
\item gravitational encounters,
\item secular perturbations,
\item orbital crossing,
\item long-term stability.
\end{itemize}
Different choices give somewhat different notions of "clearing."
So there is no simple physical surface:
inside = planet, outside = not planet.

 The distinction is nevertheless physically useful
This is important: the criterion isn't nonsense.
Consider Earth and Pluto.
Earth has mass
$$ M_E \approx 5.97\times10^{24}\ {\rm kg}, $$
whereas Pluto has
$$ M_P \approx 1.31\times10^{22}\ {\rm kg}. $$
The difference is about a factor of 450.
Consequently Earth has had an enormous gravitational influence on its orbital neighborhood, whereas Pluto is dynamically just one member of a large population of trans-Neptunian bodies.
So dynamically there really is a profound difference between Earth and Pluto.
The controversy is whether that difference should determine the meaning of the word planet.

An even more revealing example: the asteroid belt
Consider Ceres.
Ceres is:
approximately spherical,
differentiated,
gravitationally rounded,
much larger than most asteroids.
If the criterion were simply
$$ \text{approximately spherical body orbiting the Sun}, $$
then Ceres would naturally be classified as a planet.
But Ceres sits in the asteroid belt and has not become dynamically dominant there.
Hence:
$$ \text{Ceres} \rightarrow \text{dwarf planet}. $$
This gives the classification a rather interesting structure:
$$ \begin{array}{c} \text{Mercury, Venus, Earth, Mars, ...}\\ \text{dynamically dominant}\\[4pt] \Downarrow\\ \textbf{planets} \end{array} $$
versus
$$ \begin{array}{c} \text{Ceres, Pluto, Eris, ...}\\ \text{not dynamically dominant}\\[4pt] \Downarrow\\ \textbf{dwarf planets}. \end{array} $$
The second group can nevertheless contain objects that are much more planet-like physically than many small moons.

\paragraph{ The Moon exposes another oddity}
Suppose we ask why the Moon isn't a planet.
It is certainly round and massive enough to be in hydrostatic equilibrium. But it doesn't independently orbit the Sun in the simple sense required by the IAU definition; it orbits Earth while the Earth-Moon system orbits the Sun.
Yet dynamically, the Moon is a remarkably planet-like object.
And imagine moving the Moon into an independent solar orbit.
Its physical properties wouldn't change, but its classification would.
Again, classification depends partly on orbital architecture rather than solely on physical nature.

\paragraph{ The underlying philosophical choice}
This is really the heart of the controversy.
There are two legitimate ways of organizing the Solar System.
Dynamical taxonomy:
$$ \boxed{\text{Planet} = \text{body that dominates its orbital zone}} $$
This gives the IAU's eight planets a very clean dynamical hierarchy.
Or one can use a geophysical taxonomy:
$$ \boxed{\text{Planet} = \text{self-gravitating, approximately round planetary body}} $$
Then Pluto, Ceres, Eris, and perhaps several other objects join the planetary family.
Neither description changes any physical observation. They organize the same physical objects according to different principles.

There is a particularly striking asymmetry:
The IAU definition makes "planet" partly relational.
It is not simply
$$ P=P(\text{mass, radius, composition, structure}), $$
but more like
$$ P=P(\text{body},\text{orbit},\text{neighbors},\text{dynamical history}). $$
That isn't necessarily bad—many physical classifications are relational. But it means that "planet" is not a purely intrinsic category.
And this is why the Pluto controversy is more interesting than the popular story "Pluto was demoted because it wasn't big enough." Size isn't the decisive criterion. Dynamical dominance is.
In fact, from the standpoint of planetary physics, Pluto is a fascinating and genuinely planetary-type world. The disagreement is largely over which physical distinction should be encoded in the word "planet."


\section{claude}
In 2006, astronomers decided Pluto didn't qualify as a full planet anymore — but the real story is that our understanding of the outer solar system caught up with reality.

Here's what happened: for most of the 20th century, "planet" was never formally defined — Pluto just got called one by tradition since Clyde Tombaugh found it in 1930. But starting in the 1990s, astronomers began discovering more and more icy objects way out past Neptune, in a region now called the Kuiper Belt. Pluto, it turned out, wasn't a lone oddball — it was just one of many similar bodies out there.

The tipping point came in 2005, when a team found Eris, an object in that same region that's roughly Pluto's size (actually a bit more massive, though slightly smaller in diameter). Suddenly astronomers faced a choice: either call Eris the 10th planet — and probably add several more later as other Kuiper Belt objects turned up — or draw a clearer line about what "planet" actually means.


They chose the second path. The International Astronomical Union voted in 2006 on a formal definition of "planet," requiring a body to:
\begin{itemize}
   \item Orbit the Sun.
   \item Be massive enough for gravity to pull it into a round shape.
   \item Have "cleared its orbital neighborhood" — meaning it gravitationally dominates its orbit, having swept up or ejected most other debris nearby.
\end{itemize}
Pluto passes the first two but fails the third — it shares its orbital zone with a swarm of other Kuiper Belt objects and doesn't dominate it the way Earth, Jupiter, or the other eight do. So it got reclassified into a new category: dwarf planet, alongside Eris, Ceres (previously called an asteroid), Haumea, and Makemake.

So it's less that Pluto changed, and more that we found its neighbors and had to decide what "planet" should mean once we knew they were there.

Not all astronomers agree — this is actually still a live, unresolved fight among astronomers, not settled consensus.

A few things worth knowing:

The 2006 vote itself was contentious and narrow. After a week of contentious public and private debate, a small cluster of astronomers voted to demote Pluto, with the IAU defining a planet as something that orbits the sun, is massive enough to be round, and has cleared the neighborhood around its orbit. Only a few hundred of the IAU's thousands of members were even in the room for the final vote. 
Scientific American

Prominent astronomers still push back. Alan Stern, who leads NASA's New Horizons mission that flew past Pluto, has argued for years that the "clearing the neighborhood" rule is flawed and that Pluto should still count as a planet — while Mike Brown, the astronomer whose discovery of Eris triggered the whole reclassification debate, still defends the demotion. Many leading planetary scientists still disagree about how planets should be defined and whether Pluto deserves to regain full planet status. 

There's a rival "geophysical" definition on the table. Some scientists have proposed a simpler, more inclusive definition: a planet is any round object in space smaller than a star. Under that rule, counting up a little over 100 objects that would probably qualify, with more likely to be discovered — including Pluto, but also things like Europa, Titan, and Earth's own Moon. 

It's even become political. As of earlier this year, NASA's administrator, speaking to a congressional committee, said he was very much in the camp of "make Pluto a planet again" and that NASA was taking steps toward that. 

So the "8 planets" answer is really the official IAU position, not a unanimous scientific one — plenty of working astronomers and planetary scientists still call Pluto a planet, and the argument has run hot for essentially the whole 20 years since the demotion.

\section{gemini}

Pluto was reclassified in 2006 from a full planet to a dwarf planet. Nothing changed about Pluto itself—astronomers simply changed the official scientific definition of a "planet."

Why Pluto Was Reclassified
When Pluto was discovered in 1930, astronomers thought it was much larger than it actually is. By the 1990s and early 2000s, better telescopes revealed thousands of icy objects past Neptune in a region known as the Kuiper Belt.

In 2005, astronomers discovered Eris, an object in the Kuiper Belt that was roughly the same size as Pluto. This forced scientists to make a choice: either make Eris (and potentially dozens of similar icy bodies) the 10th, 11th, and 12th planets, or officially define what a "planet" actually is.

The 3 Rules of a Planet
In 2006, the International Astronomical Union (IAU) established three criteria for a body to be considered a major planet:
\begin{enumerate}
   \item It must orbit the Sun. (Pluto passes this)
   \item It must have enough mass for its gravity to pull it into a round shape. (Pluto passes this)
   \item It must have "cleared its neighborhood" around its orbit. (Pluto fails this)
\end{enumerate}
Pluto failed the third rule because its orbit is cluttered with thousands of other icy Kuiper Belt objects, and its mass isn't large enough to gravitational dominate or clear away those surrounding objects. As a result, the IAU created a new category called dwarf planets, with Pluto as the flagship member alongside Eris, Ceres, Haumea, and Makemake.
 
 \subsection{ Ceres (The Inner Solar System Outlier)}
Location: Main Asteroid Belt (between Mars and Jupiter)

Key Feature: Ceres is the only dwarf planet located in the inner solar system; all others lie far beyond Neptune.

Details: Discovered in 1801, Ceres was originally considered a planet for decades before being reclassified as an asteroid, and then finally as a dwarf planet in 2006. It accounts for about a third of the entire mass of the asteroid belt. NASA’s Dawn spacecraft visited Ceres in 2015, revealing bright salt deposits and evidence of frozen water beneath its surface.

\subsection{ Eris (The Planet Killer)}

Location: Scattered Disc (beyond the Kuiper Belt)

Key Feature: The discovery of Eris directly triggered Pluto's demotion.

Details: Discovered in 2005, Eris is almost identical in size to Pluto (about 27\% more massive, though slightly smaller in volume). Because it is so far out—taking 557 years to orbit the Sun—its surface is extremely cold and covered in frozen methane ice. It has one moon, named Dysnomia.

\subsection{Haumea (The Fast Spinner)}

Location: Kuiper Belt

Key Feature: Shaped like an elongated football or egg rather than a sphere.

Details: Haumea rotates on its axis once every 4 hours, making it one of the fastest-rotating large objects in the solar system. This rapid rotation stretched the body into an oval shape. In 2017, astronomers discovered that Haumea has a ring system, making it the first dwarf planet known to possess rings. It also has two moons, Hiʻiaka and Namaka.

\subsection{ Makemake (The Red Giant of the Kuiper Belt)}
Location: Kuiper Belt

Key Feature: Exceptionally bright with a reddish hue.

Details: Discovered shortly after Easter in 2005 (and named after the Rapa Nui creation deity of Easter Island), Makemake is the second-brightest object in the Kuiper Belt after Pluto. Its surface is covered in frozen methane and ethane. In 2016, astronomers discovered a small, dark moon orbiting Makemake, nicknamed "MK2."
\begin{tabular}{ccccc}
    Dwarf Planet&	Region&	Orbital Period	&Diameter &	Moons\\
    Ceres &	Asteroid Belt&	4.6 years&	950 km &	0\\
    Pluto	&Kuiper Belt&	248 years	&2,377 km &	5\\
    Haumea	&Kuiper Belt&	285 years&	1,600 km &	2\\
    Makemake&	Kuiper Belt&	305 years&	1,430 km &	1\\
    Eris&	Scattered Disc&	557 years&	2,326 km &	1
\end{tabular}
(Note: Astronomers estimate there may be hundreds of additional dwarf planets in the Kuiper Belt waiting to be officially designated, including prominent candidates like Gonggong, Quaoar, and Sedna.)


The reclassification of Pluto by the International Astronomical Union (IAU) in 2006 remains one of the most heavily debated decisions in modern science. Many planetary scientists and astronomers, including key leaders from NASA’s New Horizons mission, strongly reject the IAU's definition and continue to treat Pluto as a true planet.

\subsection{The "Clearing the Orbit" Rule Is Flawed and Context-Dependent}
Critics argue that the IAU’s third requirement—that a body must "clear its neighborhood" around its orbit—is flawed because it defines a planet by where it is located rather than what it is.

Location bias: If Earth were placed in Pluto’s orbit, it would not have enough gravitational dominance to clear out the massive volume of space in the Kuiper Belt either. Critics ask: Does a world stop being a planet simply because you move it farther away from the Sun?

No planet has fully cleared its orbit: Jupiter shares its orbit with over 10,000 Trojan asteroids, Earth regularly crosses paths with near-Earth asteroids, and Neptune hasn't cleared Pluto itself.

\subsection{ Geological Complexity Should Define a Planet}
Planetary scientists tend to favor a geophysical definition: A planet is a sub-stellar mass body that has never undergone nuclear fusion and has sufficient self-gravitation to assume a spherical shape.

When NASA's New Horizons spacecraft flew past Pluto in 2015, it revealed a world far more complex than many expected:
\begin{itemize}
   \item 
High nitrogen ice mountains, glaciers, and active ice volcanoes.
   \item 
A thin atmosphere with multiple haze layers and weather patterns.
   \item 
A subsurface liquid water ocean beneath its crust.
   \item 
Five distinct moons, including Charon, which shares a complex barycenter with Pluto.
\end{itemize}
To planetary scientists, a body with active geology, atmospheres, and moons functions like a planet in every meaningful scientific way.


In scientific literature, planetary scientists routinely use the word "planet" to describe worlds based on their physical properties rather than IAU rules. Under the geophysical definition, bodies like Pluto, Ceres, Europa, and Titan are all studied as planets.


Opponents also point out procedural issues with the 2006 vote itself. Out of nearly 10,000 IAU members worldwide, fewer than 500 astronomers were present on the final day of the conference in Prague to cast a vote. Furthermore, the majority of attendees were dynamicists (who study orbits) rather than planetary geologists (who study the worlds themselves).

